%\input ../util/bookmacros
%\input ../util/docparams

\appendix{Appendix F}{How to Study from This Book}
This book was written to give you a map of calculus: the four strategic
moves, the discrete and continuous boards, and the dictionary that ties
them together. But a map you never unfold fades. This appendix is about
keeping the map alive after you close the book. It rests on a small
number of findings about human memory -- findings which are as well
established as anything in the study of learning\footnote{\kern 0.5pt
\raise 0.5ex \hbox{\dag}}{They come from a century of experiments,
starting with Ebbinghaus's measurements of his own forgetting in the
1880s, through the modern work on ``retrieval practice'' by Roediger and
Karpicke (2006), the ``spacing'' and ``desirable difficulties'' work of
the Bjorks, and Skemp's 1976 comparison of relational and instrumental
understanding. The findings have held up; education fashion comes and
goes around them.} -- and on a piece of good luck: you have something no
student in history has had, a tireless examiner. Used correctly, an AI
will keep this material alive in you for years. Used incorrectly, it
will do your thinking for you and you will retain nothing. The
difference between the two is the subject of this appendix.

\section{Four Findings About Memory}
Just as the calculus in this book was organized around four strategic
moves, the practice of {\it keeping\/} it is organized around four
findings about memory.

\beginEnum
\enum{{\it Retrieval beats rereading.} Rereading a chapter feels like
learning -- the material seems familiar, and familiarity masquerades as
knowledge. But the act that strengthens memory is pulling the idea {\it
out\/} of your head, not putting it in again. Testing yourself is not
merely a way to measure knowledge; it is the way to create it.}
\enum{{\it Spaced beats massed.} Ten problems tonight do less for you
than three tonight, three next week, and three next month. Memory is
strengthened most when you retrieve something just as you are on the
verge of losing it. Forgetting a little, then recovering, is not a
failure of the method -- it {\it is\/} the method.}
\enum{{\it Mixed beats blocked.} Twenty integration-by-parts problems in
a row teach you to execute integration by parts. They do not teach you
to {\it recognize\/} when it applies -- because in a block of twenty,
recognition is never in question. Real problems arrive unlabeled. Your
practice should too.}
\enum{{\it Struggle first, then check.} An answer you fought for -- even
if you lost -- sticks far better than an answer you were shown. Commit
to an answer, in writing, {\it before\/} you see the solution. The
discomfort of not knowing is not an obstacle to learning; it is the
sensation of learning.}
\endEnum

Notice what these findings say about the traditional arrangement --
read a large book slowly, do a block of like problems after each
section, review everything the week of the exam. It is almost perfectly
backwards. It produces what it optimizes for: a few weeks of buffered
material, released on exam day and gone by summer.

\section{The Method}
Here is the whole method. For each chapter of this book:

\beginEnum
\enum{{\it Read the chapter}, working the examples with pen and paper as
you go. This is the only step that resembles ordinary studying.}
\enum{{\it Close the book and rebuild the summary.} Every chapter ends
with a summary. Before you reread it, write your own from memory: the
definitions, the theorems, the moves that were played. Then compare.
What you missed is not a mark against you -- it is the most valuable
information the exercise produces, because it tells you exactly what to
revisit.}
\enum{{\it Work the exercises cold.} Commit to an answer before
checking anything. If you are stuck for ten minutes, that is fine;
finding 4 above says the struggle is doing the work.}
\enum{{\it Then bring in the examiner} -- the next day, the next week,
and the next month, as laid out below.}
\endEnum

The whole apparatus costs perhaps two hours a week. It is not the
grind; it is the thing the grind was always pretending to be.

\section{Rules for the Examiner}
An AI can play two roles: oracle or examiner. As an oracle it answers
your questions, solves your problems, and explains what you ask it to
explain. This feels wonderful and teaches you almost nothing -- it is
rereading, with better production values. As an examiner it asks {\it
you\/} the questions, waits for your committed answer, and only then
tells you the truth. Everything good in this appendix comes from
holding the AI in the examiner's chair. Three rules accomplish this:

\beginEnum
\enum{{\it Never let it solve before you commit.} Tell it so, in the
prompt, explicitly: ``Do not show me the answer or any hint until I have
given my answer.'' A good examiner is told the rules of the exam.}
\enum{{\it Make it grade the setup separately from the algebra.} In
this book the hard step was never the arithmetic; it was choosing the
move -- recognizing the telescoping sum, seeing that an integrating
factor is called for. Instruct the examiner to grade ``did you choose
the right move'' and ``did you execute it'' as two separate marks. An AI
will happily absorb the arithmetic burden anyway; what must remain
yours is the choice of move.}
\enum{{\it Make it probe for understanding, not recall.} The test of
understanding is not whether you remember a formula -- it is whether
you can {\it rebuild\/} it from the move that produced it, and whether
you know what it is {\it for\/}. Years ago I watched a room of bright
engineering students fail an exam problem with the formula for
integration by parts written on the board in front of them. They had
memorized everything and understood what none of it was for. The
examiner's best questions are the ones that would have exposed this a
year earlier: ``What does this formula trade for what?'' ``When would
you reach for it?'' ``Derive it from the product rule.''}
\endEnum

\section{Prompts That Work}
What follows are prompts you can give an AI, nearly verbatim. They are
grouped by chapter, and each one embodies the rules above. You will
naturally drift from these scripts as you and your examiner develop a
working relationship; that is as it should be. What should not drift is
the structure: examiner not oracle, commit before checking, setup
graded apart from algebra.

\subsection{After the chapters on Difference and Summation Calculus}
{\medskip\narrower\noindent\raggedright {\tt You are my examiner on difference
calculus. Give me sequences, one at a time, and I will compute their
difference sequences. Then reverse the game: give me difference
sequences and I will find antidifferences. Mix in sums I should
evaluate by recognizing a telescoping difference and applying the
Fundamental Theorem of Difference Calculus. Do not show answers or
hints until I commit to an answer. Keep track of the kinds of errors I
make and return to those kinds later without telling me you are doing
so.}\par\medskip}

The last sentence quietly gives you interleaving and spacing inside a
single session: the examiner reintroduces your weak spots when you are
no longer expecting them.

\subsection{After the chapters on Differential and Integral Calculus}
{\medskip\narrower\noindent\raggedright {\tt Quiz me on differentiation and
integration, mixed together and unlabeled, one problem at a time. After
each of my answers, grade two things separately: whether I chose the
right technique (linearity, product rule, chain rule, substitution,
integration by parts) and whether I executed it correctly. Do not tell
me which technique a problem needs -- recognizing it is the point.}%
\par\medskip}

And a second game, which no textbook can play with you:

{\medskip\narrower\noindent\raggedright {\tt Show me a worked evaluation of an
integral that contains exactly one subtle error. I will find the error
and explain why it is wrong. Start easy and get harder as I succeed.}%
\par\medskip}

Finding a planted error exercises precisely the verification skill you
will need in the age of AI: the machine will do the computation, and
{\it you\/} must be the one who can tell whether it is right.

\subsection{After the chapter on the Derivative as Approximation}
{\medskip\narrower\noindent\raggedright {\tt Give me quantities to estimate with
the tangent line approximation -- no calculator -- and then have me
bound the error using the next term of Taylor's Theorem. Then ask me
conceptual questions: why is the tangent line the best line? What does
the error term buy us? Make me derive the first two terms of Taylor's
Theorem from the Mean Value Theorem.}\par\medskip}

\subsection{After the chapters on Difference and Differential Equations}
{\medskip\narrower\noindent\raggedright {\tt Give me word problems -- loans,
populations, cooling coffee, radioactive decay -- without telling me
whether they are discrete or continuous. I will set up the difference
or differential equation, name the strategic move (summing factor or
integrating factor), and solve it. Grade the setup, the move, and the
algebra as three separate marks. Occasionally give me a problem where
the model breaks down and ask me where and why.}\par\medskip}

\subsection{The Monthly Examination}
Once a month, hold a real examination -- oral, in the old style:

{\medskip\narrower\noindent\raggedright {\tt You are my oral examiner for twenty
minutes on the whole of discrete and continuous calculus. Ask me to:
derive the Fundamental Theorem of Calculus from the telescoping sum
picture; explain what integration by parts trades for what, and derive
it from the product rule; state the discrete/continuous dictionary from
memory; solve one difference equation and one differential equation and
point out that they are the same game. Push back on vague answers.
At the end, tell me which parts of my understanding are load-bearing
and which have gone soft.}\par\medskip}

The one-month spacing is chosen on purpose: by finding 2, you want to
retrieve this material roughly when you are starting to lose it. If the
viva goes badly in places, that is the system working -- you have found
the soft spots while they are still cheap to repair.

\section{The Drill Companion}
This book deliberately contains few exercises; it is a map, not a
gymnasium. For raw volume of practice, pair it with any inexpensive
problem collection -- the Schaum's Outline of Calculus is the
traditional choice and costs less than lunch.\footnote{\kern 0.5pt
\raise 0.5ex \hbox{\dag}}{Any edition will do. The problems have not
changed in fifty years because the calculus has not changed in three
hundred.} Two rules keep the drill useful:

\item{Do not work the problems in blocks. Open two chapters at once and
alternate, or have your examiner select page numbers for you at random.
Finding 3 again: real problems arrive unlabeled.}
\item{Before each problem, say -- out loud or on paper -- which
strategic move it will need. {\it Then\/} work it. This one habit
converts a thousand grind-problems into a thousand recognition drills,
which is the skill the grind never taught.}

\section{A Schedule}
Collecting all of the above into one place. After finishing a chapter:

\bigskip
\centerline{\vbox{\halign{\quad\hfil#\hfil\quad&\quad#\hfil\quad\cr
{\bf When}&{\bf What}\cr
\noalign{\smallskip\hrule\smallskip}
same day&rebuild the chapter summary from memory; exercises, cold\cr
next day&examiner session on the new chapter\cr
next week&mixed session: new chapter plus everything before it\cr
each month&the twenty-minute oral examination, whole book\cr
}}}
\bigskip

\noindent Miss a session, pick up where you left off; the schedule is a
rhythm, not a debt. And notice what is absent from it: there is no
``review week,'' because there is nothing left to cram -- the material
never had a chance to leave.

\section{The Test}
A year from now, sit down with a blank sheet of paper and try to write
out the dictionary that closed the last chapter -- $\Delta$ and
$d/dt$, summation and integration, the two Fundamental Theorems, the
summing factor and the integrating factor, $a = 1 + r$ tying geometric
to exponential growth. Then pick one row and rebuild it: derive the
continuous entry from the discrete one by the {\it approximate, then
take a limit\/} move.

If you can do that, you have what this book was written to give you.
Not a buffer of formulas -- those were never the point, and the machine
holds them better than you ever will -- but the map: which game is
being played, which move applies, and the ability to rebuild any tool
you have lost. That is what ``knowing calculus'' can still mean, and
it is the part no examiner, human or artificial, can hold for you.
